\section{Conclusion}
The shared SMHI-to-ROOT workflow produced three views of the temperature record. Falsterbo's calendar-day averages show the seasonal cycle and the variability around it. The 30 June histogram resolves that variability into a distribution for a specific date and illustrates how empirical and model-based probabilities can be computed from the same yearly means. In Lund, the annual count of days with mean temperature below freezing decreases by approximately 1.32 days per decade in the archived reference series.

The principal limitation is uneven observation coverage. Changes in measurement schedules and gaps affect the daily means and can reduce annual cold-day counts. The cleaning logs and archived figures make the data selection and resulting summaries traceable. The largest remaining improvement for the Lund result would be to compare the trend across years with similar coverage during the cold months.

\section*{Use of AI tools}
\addcontentsline{toc}{section}{Use of AI tools}
AI assistance was used during development to debug the calendar-day analysis, translate familiar Python operations into C++ and diagnose compilation issues. It also supported the preparation and integration of the cleaning and cold-day modules. For this report, AI assisted in organising and drafting the LaTeX, checking descriptions against the source code and recomputing the Lund fit diagnostics from the archived annual counts. These uses supported implementation and documentation. The data, plots and analysis methods are traceable to the project repository.
